\documentclass[cameraready]{Interspeech}
\usepackage{float}

\title{Image Classification with Convolutional Neural Networks on CINIC-10}

\author[affiliation={1}]{Pawel}{Pozorski}
\author[affiliation={1}]{Pawel}{Florek}

\address{
    $^1$ Warsaw University of Technology, Poland
}

\email{\{pawel.pozorski, pawel.florek\}.stud@pw.edu.pl}

\keywords{image classification, convolutional neural networks, CINIC-10, neural architecture search, few-shot learning}

\begin{document}

\maketitle

\begin{abstract}
This paper presents a comparative study of convolutional neural network pipelines for image classification on CINIC-10. We evaluate lightweight and high-capacity backbones, including MobileNetV3 Small, SqueezeNet, ResNet-18, DenseNet-121, and a differentiable NAS-CNN variant. We also test ConvKAN layers as convolutional replacements in efficient models. The study includes a controlled hyperparameter search, modern augmentation methods, reduced-data training, and few-shot metric learning with Prototypical Networks. Experiments are designed for reproducibility with fixed seeds and unified training scripts. We report benchmark context, evaluation protocol, and implementation details that support transparent comparison across model families and training regimes.
\end{abstract}

\section{Introduction}
Convolutional neural networks (CNNs) \cite{lecun1998gradient} remain a practical foundation for image recognition \cite{liu2022convnet}, but model selection is highly task- and constraint-dependent \cite{bianco2018benchmark, tan2019efficientnet}. In this project, we study how different CNN paradigms perform on CINIC-10 \cite{cinic}, a dataset that preserves CIFAR-like \cite{krizhevsky2009learning} image size while increasing distributional complexity via ImageNet-derived \cite{deng2009imagenet} samples.

Compared with CIFAR-10, CINIC-10 introduces a stronger distribution shift because class semantics are preserved while image statistics are broadened by downsampled ImageNet content \cite{cinic}. This makes the benchmark less prone to narrow in-distribution overfitting and more informative for robustness-oriented evaluation. In practice, models that perform similarly on CIFAR-like data may diverge on CINIC-10 once they are exposed to larger intra-class variation, background diversity, and texture/illumination changes \cite{cinic, recht2019do}.

Our goals are threefold: (i) compare efficient and high-capacity architectures under a common training framework, (ii) evaluate whether ConvKAN-style convolutions improve compact models, and (iii) analyze robustness in low-data conditions using subset training and few-shot learning.

\section{Methodology}
\subsection{Dataset}
CINIC-10 \cite{cinic} contains $270{,}000$ RGB images at $32 \times 32$ resolution and $10$ balanced classes \cite{cinic}. samples from the dataset are illustrated in Figure \ref{fig:sample_img}. Data is organized into train, validation, and test splits of $90{,}000$ samples each. The inclusion of ImageNet-derived examples increases intra-class variation relative to CIFAR-10 and makes generalization analysis more informative \cite{cinic}. For the purpose of these experiments, the dataset was sourced from the \texttt{mengcius/cinic10} repository via the Kaggle API, ensuring a standardized pipeline for data loading and preprocessing.

\begin{figure}
    \centering
    \includegraphics[width=\linewidth]{../notebooks/plots/eda_sample_images.png}
    \caption{Example images from CINIC-10 for four classes.}
    \label{fig:sample_img}
\end{figure}



\subsection{Model families}
We evaluate three groups:
\begin{itemize}
    \item \textbf{Search-based architecture}: a differentiable one-shot NAS-CNN supernet with architecture parameters optimized jointly with weights \cite{SalmaniPourAvval2025_nas}.
    \item \textbf{Efficient backbones}: MobileNetV3 Small \cite{Howard_2019_ICCV} and SqueezeNet \cite{SqueezeNet}, including ConvKAN substitutions \cite{10681070_kan}.
    \item \textbf{Transfer baselines}: ResNet-18 \cite{He_2016_CVPR_res} and DenseNet-121 \cite{Huang_2017_CVPR_dense} adapted for the CINIC-10 label space.
\end{itemize}

\subsubsection{Convolutional Kolmogorov-Arnold Networks (KANs)}
For ConvKAN variants, we replace eligible spatial convolution layers with \texttt{torchkan}-provided \cite{torchkan} ConvKAN layers while preserving layer shape and stride constraints. Conceptually, this swaps standard linear convolutional kernels for KAN-style parameterization with learnable nonlinear basis behavior on connections, motivated by recent work on Kolmogorov-Arnold Networks \cite{liu2024kan,10681070_kan}. Unlike standard convolutions, which apply fixed scalar weights followed by a static, node-wise activation function (e.g., $y = \text{tanh}(\sum w_i x_i)$), KANs shift the nonlinearity directly to the network edges. In a ConvKAN layer, the scalar elements of a traditional convolutional kernel are replaced by learnable 1D nonlinear functions, typically parameterized as B-splines.  Consequently, each connection applies a dynamic, finely-tuned activation $\phi(x)$ to the input pixels within the receptive field, and the output node simply sums these transformed signals ($y = \sum \phi_i(x_i)$). This theoretical increase in per-parameter expressivity aims to allow lightweight architectures to capture complex spatial patterns more efficiently than standard linear filters. CNNs with ConvKAN layers have shown promise in benchmarks \cite{drokin2024kolmogorov}, however their application is still relatively novel and underexplored, especially in the context of efficient vision models.

\subsubsection{Neural Architecture Search (NAS)}

For NAS, our supernet is constructed as a sequence of searchable edges, formulated via a continuous relaxation of the discrete architecture space \cite{liu2018darts}.  Instead of selecting a single operation, the output of each edge during the search phase is computed as a weighted sum of all candidate operations $o \in \mathcal{O}$, combined using a softmax distribution over learnable architecture parameters $\alpha$:
$$ \bar{o}(x) = \sum_{o \in \mathcal{O}} \frac{\exp(\alpha_o)}{\sum_{o'} \exp(\alpha_{o'})} o(x) $$

Our specific search space $\mathcal{O}$ consists of $5$ candidate operations designed to balance capacity and efficiency: standard convolutions (\texttt{conv3x3}, \texttt{conv5x5}), parameter-efficient \texttt{depthwise3x3}, spatial reduction \texttt{maxpool3x3\_proj}, and identity mapping \texttt{skip\_proj}. The \texttt{\_proj} variants explicitly incorporate a $1 \times 1$ pointwise convolution to ensure tensor dimensions match when spatial strides or channel counts shift across an edge. \texttt{Maxpool3x3\_proj} applies a $3 \times 3$ max pooling followed by a $1 \times 1$ convolution to adjust the channel dimension. Each \texttt{conv} operation uses a \texttt{BatchNorm} and \texttt{SiLU} activation.

A depthwise convolution does not share weights across channels - it learns a distinct spatial filter for each input channel. For example, given an input with $32$ channels, a standard $3\times3$ convolution computing $32$ output channels requires $32 \times 32 \times 3 \times 3 = 9{,}216$ parameters to mix all spatial and channel information simultaneously. A depthwise $3\times3$ layer requires only $32 \times 3 \times 3 = 288$ parameters, as it applies exactly one unique $3\times3$ filter to each input channel independently.

After the feature extraction backbone, we apply \texttt{AdaptiveAvgPool2d}$(1,1)$ to pool feature maps from $[C=128, H, W]$ to $[128, 1, 1]$. This pooling reduces spatial dependence and the number of hyperparameters in the classifier head, which consists of: \texttt{Flatten}$(128)$ $\rightarrow$ \texttt{Dropout} $\rightarrow$ \texttt{Linear}$(128, \text{num\_classes})$. During NAS training, temperature annealing is applied to the architecture softmax, using linear annealing from $T_{\text{start}} = 5.0$ to $T_{\text{end}} = 0.5$ across epochs:
\[
T(e) = T_{\text{start}} + \frac{e-1}{E-1} \left(T_{\text{end}} - T_{\text{start}}\right),
\]
where $e \in \{1, \ldots, E\}$ is the current epoch. At high temperatures the operation mixture is exploratory; as temperature decreases, the architecture distribution sharpens toward the best operations.

The macro-architecture skeleton restricts the search to $6$ sequential edges, following a standard hierarchical feature pyramid defined by these input/output channel ($C_{in} \rightarrow C_{out}$) and stride ($s$) transitions:
\begin{itemize}
    \item $(3\rightarrow32, s=1)$, $(32\rightarrow32, s=1)$, $(32\rightarrow64, s=2)$,
    \item $(64\rightarrow64, s=1)$, $(64\rightarrow128, s=2)$, $(128\rightarrow128, s=1)$.
\end{itemize}
This 6-edge macro-skeleton was chosen empirically by the authors. This formulation enables the joint gradient-based optimization of both the network weights and the architecture parameters $\alpha$. Following the search stage, the continuous edge mixtures are discretized by taking the $\arg\max$ of the architecture logits to derive the final feature transformation and downsampling pattern.

\subsection{Training and search setup}
A structured grid search is used for MobileNetV3 Small with the Cartesian product of:
\begin{itemize}
    \item optimizer $\in \{$SGD \cite{bottou2010large}, AdamW \cite{adamw}$\}$,
    \item batch size $\in \{128, 256\}$,
    \item dropout $\in \{0.0, 0.1, 0.5\}$,
    \item weight decay $\in \{0.0, 10^{-4}\}$.
\end{itemize}

These bounds were selected to bracket key training regimes: batch sizes of $128$ and $256$ balance computational efficiency with the implicit regularization of stochastic mini-batch noise \cite{masters2018revisiting}, while the dropout values evaluate whether a lightweight architecture requires no explicit regularization ($0.0$), mild feature decorrelation ($0.1$), or classic aggressive dropout ($0.5$) \cite{srivastava2014dropout}. Choosing correct optimizer can improve computer vision tasks accuracy, with SGD and Adam variants being one of the most reliable options \cite{Hassan2023}. AdamW \cite{adamw} is included as a modern variant over Adam improving generalization by decoupling weight decay from the adaptive learning rate mechanism. Weight decay values of $0.0$ and $10^{-4}$ test the necessity of explicit regularization, which can be crucial for preventing overfitting \cite{NEURIPS2024_29496c94}.
This yields $24$ configurations per seed in Stage 1 (refer to Section~\ref{sec:execution_plan} for details). This stage is repeated for three fixed seeds $\{0, 42, 3407\}$, resulting in $72$ total runs; thus, seed is treated as an experimental repeatability axis (for stability statistics), not as a tuned hyperparameter. Learning rate is not part of the grid and is set to a value of $0.0003$, with scheduling via cosine annealing \cite{loshchilov2016sgdr} over the fixed number of epochs.

For NAS optimization, entropy regularization is also applied to architecture logits during the search phase to encourage sharper operation selection and reduce redundancy in the learned operator weights.

\subsection{Reproducibility and seed control}
The same seed policy is applied consistently across all stages in the execution plan (grid search, augmentation ablation, architecture comparison, NAS two-stage, reduced-data training, and few-shot episodes). One top-level seed initializes Python, NumPy \cite{harris2020array}, and PyTorch \cite{NEURIPS2019_9015} RNG states at process start. Dataloader randomness is controlled with split-specific seeded generators (train/validation/test offsets), so ordering is deterministic within a run but decorrelated across splits. In multi-worker loading, each worker receives a deterministic derived seed and forwards it to Python and NumPy RNGs inside that worker process.

This seed is also propagated to class-stratified subset sampling and few-shot episode construction, so data selection, task generation, augmentation sampling, and batch shuffling are all reproducible for a given run directory (e.g., \texttt{seed\_0}, \texttt{seed\_42}, \texttt{seed\_3407}). Reporting across these seeds provides mean/variance-style stability evidence rather than a single-point estimate.

\subsection{Data augmentation}
The training pipeline includes geometric and photometric transforms and mixing-based regularization:
\begin{itemize}
    \item Random rotation, horizontal flip, color jitter, and AutoAugment (CIFAR-10 policy) \cite{Cubuk_2019_CVPR} are applied inside the dataset transform pipeline (torchvision transforms) used by the DataLoader. Therefore, these operations are performed online at sample level when worker processes read images.
    \item MixUp \cite{zhang2018mixupempiricalriskminimization} and CutMix \cite{Yun_2019_ICCV_cutmix} are applied at \textbf{batch level} in the DataLoader \texttt{collate\_fn} (torchvision v2 \texttt{MixUp}/\texttt{CutMix}), i.e., after a mini-batch is sampled from the dataset and collated.
    \item For MixUp, for a batch permutation $\pi$ and $\lambda \sim \mathrm{Beta}(\alpha,\alpha)$, mixed pairs are formed as
    \[
    \tilde{x}_n = \lambda x_n + (1-\lambda)x_{\pi(n)}, \qquad
    \tilde{y}_n = \lambda y_n + (1-\lambda)y_{\pi(n)}.
    \]
    For CutMix, with binary patch mask $M_n \in \{0,1\}^{H \times W}$,
    \[
    \tilde{x}_n = M_n \odot x_n + (1-M_n) \odot x_{\pi(n)},
    \]
    and the effective label weight is area-based,
    \[
    \lambda'_n = \frac{1}{HW}\sum_{h,w} M_{n,h,w}, \qquad
    \tilde{y}_n = \lambda'_n y_n + (1-\lambda'_n)y_{\pi(n)}.
    \]
    We train with cross-entropy for soft targets:
    \[
    \mathcal{L} = -\frac{1}{B}\sum_{n=1}^{B}\sum_{c=1}^{C}\tilde{y}_{n,c}\log p_{n,c},
    \]
    where $p_{n,c}=[\mathrm{softmax}(f_\theta(\tilde{x}_n))]_c$. This is equivalent to weighted CE on paired labels (MixUp/CutMix-specific loss).
    \item All images are always normalized per channel (train/validation/test) with fixed CINIC-10 statistics: mean $(0.47889522,\ 0.47227842,\ 0.43047404)$ and standard deviation $(0.24205776,\ 0.23828046,\ 0.25874835)$.
\end{itemize}

\begin{figure}
    \centering
    \includegraphics[width=\linewidth]{../notebooks/plots/eda_augmentation_sample.png}
    \caption{Showcase of basic augmentation methods and autoaugment.}
    \label{fig:basic_aug}
\end{figure}

\begin{figure}
    \centering
    \includegraphics[width=\linewidth]{../notebooks/plots/eda_mixing_methods.png}
    \caption{Showcase of mixing methods.}
    \label{fig:aug_mixing}
\end{figure}



\subsection{Low-data evaluation}
Two complementary low-resource protocols are used:
\begin{itemize}
    \item \textbf{Subset training}: supervised training on $5\%$ of data per class, sampled with class-stratified selection. For each class, indices are sampled independently without replacement using the experiment seed, with at least one sample retained per class; this preserves class coverage and prevents label-frequency drift that appears in naive global subsampling.
    \item \textbf{Few-shot learning}: episodic training with Prototypical Networks \cite{NIPS2017_cb8da676}. To optimize the model for transferable class separation rather than closed-set classification, training proceeds through $N$-way, $K$-shot episodes. Each episode is constructed and evaluated through the following steps:
    \begin{enumerate}
        \item \textbf{Task sampling}: A support set $S$ is formed by randomly sampling $N$ classes from the training split and $K$ images per class. A disjoint query set $Q$ is then sampled from those exact same $N$ classes to serve as the episodic test set.
        \item \textbf{Latent embedding and prototype calculation}: The CNN backbone, parameterized by $\theta$, acts as an embedding function $f_\theta$ that maps images into a continuous latent space. For each class $c$, a prototype vector $\mathbf{v}_c$ is computed as the mean of its $K$ embedded support examples:
        $$ \mathbf{v}_c = \frac{1}{K} \sum_{(x_i, y_i) \in S_c} f_\theta(x_i) $$
        \item \textbf{Distance formulation}: Query images $x_q \in Q$ are passed through the same backbone $f_\theta$. The network evaluates the squared Euclidean distance $d$ between the query embedding $f_\theta(x_q)$ and each class prototype $\mathbf{v}_c$.
        \item \textbf{Prediction and optimization}: A softmax function over the negative Euclidean distances yields a probability distribution over the $N$ episodic classes:
        $$ p(y=c|x_q) = \frac{\exp(-d(f_\theta(x_q), \mathbf{v}_c))}{\sum_{c'} \exp(-d(f_\theta(x_q), \mathbf{v}_{c'}))} $$
        The model weights $\theta$ are updated via backpropagation by minimizing the negative log-probability (cross-entropy loss) of the true query labels.
    \end{enumerate}
    By continually optimizing this distance-based metric across varying episodic tasks, the network learns a generalized feature space where intra-class samples cluster tightly, making it highly effective for low-data adaptation on unseen data distributions.
\end{itemize}

\section{Related Work and Benchmark Context}
\subsection{Efficient CNN architectures}
Lightweight vision models evolved from compact macro-architectures (e.g., SqueezeNet) toward systematic use of depthwise-separable and inverted bottleneck designs in MobileNet families \cite{SqueezeNet,sandler2018mobilenetv2,Howard_2019_ICCV}. Subsequent compound scaling strategies, such as EfficientNet, further formalized the accuracy-efficiency trade-off across depth, width, and resolution \cite{tan2019efficientnet}. Our efficient-model branch is positioned in this line of work and tests whether ConvKAN substitutions can improve compact backbones under fixed training budgets.

\subsection{Differentiable NAS}
Differentiable NAS methods treat architecture choice as a continuous relaxation optimized by gradient descent, as popularized by DARTS \cite{liu2018darts}. Hardware-aware variants such as ProxylessNAS further emphasize practical deployment constraints \cite{cai2018proxylessnas}. Our NAS-CNN setup follows this paradigm by optimizing operation-mixture logits and then retraining a discretized architecture selected from the learned supernet.

\subsection{Few-shot learning context}
Few-shot learning approaches can be broadly grouped into metric-learning and optimization-based strategies. Prototypical Networks represent the metric-learning family by classifying queries via distances to class prototypes in embedding space \cite{NIPS2017_cb8da676}, whereas MAML and related methods optimize rapid adaptation through meta-learned initialization \cite{finn2017maml}. We use ProtoNet because it integrates cleanly with the planned episodic pipeline and provides a stable low-data baseline for comparison with reduced-data supervised training.

\subsection{CINIC-10 benchmark context}
Published CINIC-10 baselines report representative validation error rates around $12.42\%$ for ResNet-18, $8.74\%$ for DenseNet-121 (train+validation protocol), and higher error for lightweight MobileNet variants \cite{cinic}. These numbers provide an external reference scale for interpreting our results.

\section{Experimental Design and Reporting}
\subsection{Experiment execution plan}
\label{sec:execution_plan}
Our experiments follow a fixed staged protocol consistent with the project execution plan.

\textbf{Stage $0$ (initial grid search).} For each seed, we run a first grid search call (refer to next stage) with $100$ epochs to identify the elbow point in the training curve where validation performance plateaus. This informs the epoch budget for subsequent stages, ensuring that all models are trained to convergence without unnecessary overfitting or undertraining.

\textbf{Stage $1$ (grid search on MobileNetV3-Small).} For each seed, we run a $24$-configuration grid over optimizer, batch size, dropout, and weight decay. Across three seeds this yields $72$ runs in total. The resulting \texttt{grid\_results.csv} files are used to select \texttt{BEST\_*} training settings.

\textbf{Stage $2$ (augmentation ablation).} Using Stage-$1$ \texttt{BEST\_*} settings, we compare $4$ augmentation configurations on MobileNetV3-Small: (i) geometric+photometric only, (ii) +MixUp, (iii) +CutMix, and (iv)
AutoAugment. The selected policy is exported as \texttt{BEST\_AUG\_EXTRA\_ARGS}.

\textbf{Stage $3$ (architecture comparison).} With fixed \texttt{BEST\_*} and best augmentation, we train SqueezeNet, ResNet-18, DenseNet-121, ConvKAN variants, and run NAS two-stage search+retrain per seed with early stopping. This stage provides the final architecture ranking under a common protocol.

\textbf{Stage $4$ (final low-data evaluation).} We run $2$ final protocols with the chosen final configuration: class-stratified reduced-data supervised training ($5\%$ per class) and few-shot episodic training (ProtoNet, $5$-way/$5$-shot/$15$-query).

\subsection{Evaluation criteria}
Primary metrics are classification accuracy and error rate on validation and test splits. For model comparison, we additionally track parameter counts, training stability across seeds, and resource usage (CPU, GPU, memory, wall-clock time). For NAS, we report both the final architecture's performance and the search phase's efficiency metrics.

\subsection{Reproducibility and compute reporting}
All experiments are executed through a unified command interface with explicit seed control. We log per-epoch (or per-episode) wall-clock and CPU time and memory usage to support reproducibility and comparison of computational cost.

Experiments are run on macOS with Apple Metal backend enabled through PyTorch MPS execution (\textit{Apple M4 16gb RAM MacBook Pro}, only default system tasks running in background, plug-in performance mode). The software stack follows project-pinned requirements: Python $\geq 3.13$, PyTorch $\geq 2.8.0$ \cite{NEURIPS2019_9015}, and torchvision $\geq 0.23.0$ \cite{torchvision}. This context allows interpreting runtime and memory numbers as platform-specific measurements rather than framework-agnostic absolutes. For more specifics, refer to the repository \footnote{\url{https://github.com/Pawlo77/cinic10}}.

\subsection{Limitations}
The complete hyperparameter grid is applied only to MobileNetV3 Small due to computational budget. Results for other families rely on reduced tuning and therefore should be interpreted as strong reference configurations rather than exhaustive optima. Our work is based only on the CINIC-10 dataset, so generalization to other benchmarks or real-world applications may be limited. The NAS search space is relatively small and may not capture all promising architectures. Finally, the few-shot evaluation focuses on Prototypical Networks, so conclusions about low-data performance may not extend to optimization-based meta-learning methods.

\section{Results}
This section will present the results of the experiments as outlined in the execution plan (refer to subsection \ref{sec:execution_plan}). We will present the findings for each stage, starting with the initial grid search, followed by the augmentation ablation, architecture comparison, NAS results, and finally the low-data evaluation. Each subsection will include relevant tables and figures to illustrate the performance of the models under different configurations and training regimes. For the sake of clarity and conciseness, we will focus on the key findings and trends observed in the results, while also discussing any notable observations or insights that emerged from the experiments. For additional details on the results, please refer to the accompanying code repository and the appended figures.

\subsection{Stage 0}

\begin{figure}[t]
    \centering
    \includegraphics[width=\linewidth]{../notebooks/plots/stage0_grid.png}
    \caption{Accuracy and loss on validation set for MobileNetV3 Small across three seeds for stage 0 initial grid search.}
    \label{fig:stage0_grid}
\end{figure}

This stage involved running the initial grid search for MobileNetV3 Small across three seeds to identify the optimal epoch budget for training. The results are summarized in Figure \ref{fig:stage0_grid}, which shows the accuracy and loss on the validation set for each seed. The optimal epoch budget is determined by identifying the point at which the validation performance plateaus, indicating that further training would not yield significant improvements and may lead to overfitting. The selected epoch budget of $60$ will be used for all subsequent stages to ensure consistent training across models and configurations.

\subsection{Stage 1}
In this stage we conducted a comprehensive grid search for MobileNetV3 Small across three seeds, evaluating the impact of optimizer choice, batch size, dropout rate, and weight decay on perfomance. For the top 20 configurations, we generated a parallel coordinates plot (Figure \ref{fig:stage1_parallel_coordinates}) to visualize the relationships between hyperparameters and test accuracy. This plot allows us to identify trends and interactions among the hyperparameters. For instance, we can observe that optimizer choice and batch size yield the best perfomance, while dropout and weight decay have smaller impact on the top configurations. To further analyze the best configurations, we created stability box plots (Figure \ref{fig:stage1_top5_stability_boxplots}) for the top 5 configurations across the three seeds.

\begin{figure}[H]
    \centering
    \includegraphics[width=\linewidth]{../notebooks/plots/stage1_parallel_coordinates_top20.png}
    \caption{Parallel coordinates plot for the top 20 hyperparameter configurations of MobileNetV3 Small.}
    \label{fig:stage1_parallel_coordinates}
\end{figure}


\begin{figure}[H]
    \centering
    \includegraphics[width=\linewidth]{../notebooks/plots/stage1_top5_stability_boxplots.png}
    \caption{Stability box plots for the top 5 hyperparameter configurations of MobileNetV3 Small.}
    \label{fig:stage1_top5_stability_boxplots}
\end{figure}

Having in mind the stability and performance, we deciced to not move with the configuration with the highest median accuracy, but with the one that has the best balance between median and stability across seeds. Therefore, we selected the best configuration for MobileNetV3 Small as follows: optimizer = AdamW, batch size = 128, dropout = 0.5, weight decay = 0.0001. This configuration will be used as the baseline for subsequent stages, including the augmentation ablation and architecture comparison. Additional analysis of training efficiency and resource usage is shown in the appendix (Figures \ref{fig:stage1_bubble_speed_accuracy_memory} and \ref{fig:stage1_violin_batchsize_vs_peak_memory}), which reveal the relationship between optimizer choice, batch size, training time, and peak memory consumption across the grid search configurations.

\subsection{Stage 2}
In this stage, we evaluated the impact of different data augmentation strategies on the performance of MobileNetV3 Small using the best hyperparameter configuration identified in Stage 1. We compared four augmentation configurations: (i) standard augmentation methods (rotations, flips, color jittering), (ii) standard + MixUp, (iii) standard + CutMix, and (iv) AutoAugment (Figure \ref{fig:stage2_augmentation_comparison}). The results indicate that applying any augmentation strategy improves performance compared to no augmentation, with standard-only augmentation method providing the best results. Interestingly, the more complex policies did not yield better performance. Standard augmentation methods will be used for all subsequent stages. Detailed analysis of augmentation strategies, including relationships between augmentation methods and training time, generalization gaps, convergence behavior, and training loss smoothness, is provided in the appendix (Figures \ref{fig:stage2_augmentation_accuracy_vs_total_time_aggregated}, \ref{fig:stage2_generalization_gap_aug}, \ref{fig:stage2_augmentation_convergence_delay_train_val}, and \ref{fig:stage2_augmentation_loss_smoothness_tail_variance}).

\begin{figure}[H]
    \centering
    \includegraphics[width=\linewidth]{../notebooks/plots/stage2_augmentation_final_stability_boxen.png}
    \caption{Comparison of different data augmentation strategies for MobileNetV3 Small.}
    \label{fig:stage2_augmentation_comparison}
\end{figure}

\subsection{Stage 3}
This stage involved a comprehensive comparison of different CNN architectures, including MobileNetV3 Small, SqueezeNet, ResNet-18, DenseNet-121, NAS discovered architecture, ConvKAN variant of MobileNetV3 Small, and ensemble model combining smaller models (MobileNetV3 Small, SqueezeNet, NAS model). Each model was trained using the best hyperparameter configuration and augmentation strategy identified in previous stages. Bigger models like ResNet-18 and DenseNet-121 achieved the best performance, while the NAS discovered architecture and ConvKAN variant of MobileNetV3 Small showed competitive results. The ensemble model, which combines predictions from multiple smaller models, also performed well, demonstrating the benefits of model ensembling in improving accuracy (Figure \ref{fig:stage3_architecture_comparison}).

\begin{figure}[H]
    \centering
    \includegraphics[width=\linewidth]{../notebooks/plots/stage3_models_mean_test_accuracy.png}
    \caption{Comparison of different CNN architectures.}
    \label{fig:stage3_architecture_comparison}
\end{figure}

Unfortunately, due to computational constraints, we were not able to run the full training on all seeds for ConvKAN variants, only $30$ epochs on seed $0$ were completed. The preliminary results suggest that ConvKAN substitutions produce similar performance to the standard MobileNetV3 Small, but with the great cost of training time, which is not a good trade-off. Wall time per epoch for ConvKAN was significantly higher than the standard MobileNetV3 Small ($> 30\times$).

The NAS search process yielded a novel architecture that performed competitively with the hand-designed models, demonstrating the potential of automated architecture search in discovering efficient and effective CNN designs for image classification tasks. It showed that the search phase is time-consuming, but the final architecture is very efficient in terms of parameter count and achieves comparable accuracy to the smaller models (Table \ref{tab:stage3_summary}). Additional analysis of model convergence and NAS architecture selection is provided in the appendix (Figures \ref{fig:stage3_accuracy_time_memory_bubble}, \ref{fig:stage3_convergence_vs_epoch_retrain_only}, and \ref{fig:stage3_nas_selected_ops_distribution}), which illustrate the training dynamics across models and the distribution of operations selected by the NAS algorithm.


\begin{table}[H]
    \centering
    \caption{Stage 3 summary: model accuracy, training time, and parameter count.}
    \label{tab:stage3_summary}
    \resizebox{\linewidth}{!}{%
    \begin{tabular}{|l|c|c|c|}
        \hline
        Model & Mean test accuracy & Mean total wall time [s] & Num. params [M] \\
        \hline
        MobileNetV3 Small & 0.5579 & 1209.28 & 1.528 \\
        SqueezeNet & 0.5864 & 690.66 & 0.741 \\
        ResNet-18 & 0.7232 & 774.64 & 11.182 \\
        DenseNet-121 & 0.7600 & 1475.67 & 6.964 \\
        NAS two-stage & 0.5513 & 17279.28 & 0.240 \\
        \hline
    \end{tabular}%
    }
\end{table}

\subsection{Stage 4}
In the final stage, we evaluated the performance of the best perfoming model from the previous stage (DenseNet-121) under low-data conditions. We conducted two complementary protocols: (i) supervised training on a class-stratified subset containing $5\%$ of the data per class, and (ii) few-shot episodic training using Prototypical Networks with a $5$-way/$5$-shot/$15$-query setup. ProtoNets showed better performance in terms of accuracy and time compared to the reduced-data supervised training, demonstrating the effectiveness of metric-based few-shot learning in low-data regimes (Figure \ref{fig:stage4_low_data_comparison}).

\begin{figure}[H]
    \centering
    \includegraphics[width=\linewidth]{../notebooks/plots/stage4_wall_time_vs_accuracy_scatter.png}
    \caption{Comparison of low-data evaluation protocols: class-stratified subset training vs. few-shot episodic training.}
    \label{fig:stage4_low_data_comparison}
\end{figure}

The results that standard DenseNet-121 trained on a reduced dataset suffered a significant drop in accuracy compared to the full-data training, while the Prototypical Network approach maintained a more robust performance despite the limited data. This highlights the importance of specialized low-data learning techniques when working with scarce labeled data, as traditional supervised training may not be sufficient to achieve good generalization in such scenarios.

\section{Conclusion}
This study defines a consistent and reproducible framework for comparing CNN architectures on CINIC-10 across efficiency, capacity, and low-data regimes. By combining NAS, ConvKAN variants, transfer baselines, and few-shot methods under shared protocols, the project provides a practical basis for selecting models under different resource and data constraints.

The results highlight the trade-offs between model complexity and performance, with larger architectures like DenseNet-121 achieving the best accuracy but at a higher computational cost. The NAS-discovered architecture showed promise in terms of parameter efficiency but required significantly more training time, while ConvKAN substitutions did not yield a favorable accuracy-time trade-off. In low-data settings, Prototypical Networks outperformed traditional supervised training, underscoring the value of metric-based approaches for few-shot learning. Tested grid search provided insights into effective training configurations with optimizer choice and batch size being particularly influential for MobileNetV3 Small. Augmentation strategies improved performance, but more complex policies did not outperform standard geometric and photometric transformations. Overall, this work contributes a comprehensive benchmark and analysis of CNN paradigms on CINIC-10, with implications for model selection in practical image classification tasks.

In future work, we suggest to expand the NAS search space and explore additional low-data learning methods. Due to computational constraints, we were not able to run the full training on all seeds for ConvKAN variants, so a more thorough evaluation of these models would be valuable. Implementing more few-shot learning approaches, such as optimization-based meta-learning methods (e.g., MAML), would provide a broader perspective on low-data performance. Additionally, testing the generalization of our findings to other datasets and real-world applications would further validate the practical relevance of our conclusions.

\section{AI Disclosure}
This report was edited with AI assistance, to improve clarity, consistency, and formatting. Part of the docstrings and some phrasing were refined using AI tools, but all scientific decisions and interpretations are the authors' own. The core content, experimental design, and analysis were developed by the authors.

\bibliographystyle{IEEEtran}
\bibliography{mybib}

\newpage
~\vfill
\newpage
\appendix
\section{Appendix}
\subsection{Additional figures}
\begin{figure}[H]
    \centering
    \includegraphics[width=\linewidth]{../notebooks/plots/stage1_bubble_speed_accuracy_memory.png}
    \caption{Bubble plot showing the relationship between test accuracy and training time for different optimizers of MobileNetV3 Small.}
    \label{fig:stage1_bubble_speed_accuracy_memory}
\end{figure}

\begin{figure}[H]
    \centering
    \includegraphics[width=\linewidth]{../notebooks/plots/stage1_violin_batchsize_vs_peak_memory.png}
    \caption{Violin plot showing the relationship between batch size and peak memory usage for MobileNetV3 Small.}
    \label{fig:stage1_violin_batchsize_vs_peak_memory}
\end{figure}

\begin{figure}[H]
    \centering
    \includegraphics[width=\linewidth]{../notebooks/plots/stage2_augmentation_accuracy_vs_total_time_aggregated.png}
    \caption{Relationship between augmentation strategies and test accuracy for different training times.}
    \label{fig:stage2_augmentation_accuracy_vs_total_time_aggregated}
\end{figure}

\begin{figure}[H]
    \centering
    \includegraphics[width=0.9\linewidth]{../notebooks/plots/stage2_generalization_gap_aug.png}
    \caption{Generalization gap for different augmentation strategies.}
    \label{fig:stage2_generalization_gap_aug}
\end{figure}

\begin{figure}[H]
    \centering
    \includegraphics[width=0.9\linewidth]{../notebooks/plots/stage2_augmentation_convergence_delay_train_val.png}
    \caption{Convergence delay for different augmentation strategies.}
    \label{fig:stage2_augmentation_convergence_delay_train_val}
\end{figure}

\begin{figure}[H]
    \centering
    \includegraphics[width=0.9\linewidth]{../notebooks/plots/stage2_augmentation_loss_smoothness_tail_variance.png}
    \caption{Loss smoothness for different augmentation strategies.}
    \label{fig:stage2_augmentation_loss_smoothness_tail_variance}
\end{figure}

\begin{figure}[H]
    \centering
    \includegraphics[width=\linewidth]{../notebooks/plots/stage3_accuracy_time_memory_bubble.png}
    \caption{Bubble plot showing the relationship between test accuracy and training time for different optimizers of MobileNetV3 Small.}
    \label{fig:stage3_accuracy_time_memory_bubble}
\end{figure}

\vfill
\newpage

\begin{figure}[H]
    \centering
    \includegraphics[width=\linewidth]{../notebooks/plots/stage3_convergence_vs_epoch_retrain_only.png}
    \caption{Convergence vs. epoch for models.}
    \label{fig:stage3_convergence_vs_epoch_retrain_only}
\end{figure}

\begin{figure}[H]
    \centering
    \includegraphics[width=\linewidth]{../notebooks/plots/stage3_nas_selected_ops_distribution.png}
    \caption{Distribution of selected operations in NAS.}
    \label{fig:stage3_nas_selected_ops_distribution}
\end{figure}


\end{document}
